\abstract{Facial sensing is often proposed as a route to affect-aware adaptation, but categorical emotion recognition is a poor default assumption for a child-facing ADHD intervention. This companion study asks a narrower question: whether locally derived, within-child facial-behavior dynamics add reproducible information about excessive task demand or recovery after task performance, explicit child input, and non-camera telemetry are already known. The module uses facial landmarks, head pose, tracking confidence, and selected temporal blendshape dynamics without assigning categorical emotion labels. Its primary evaluation is incremental validity under participant-aware out-of-sample testing. A base model contains task, context, self-report, and optional inertial features; an extended model adds facial features. Camera processing is on-device, raw video is excluded from standard telemetry, and participation in the camera module requires a separable opt-in. A facial module is retained only if it improves prespecified predictive or decision metrics robustly across participants and recording conditions. This design treats null results as a reason to simplify the adaptive controller rather than reinterpret facial signals post hoc.}
